Perch is a Soroban smart account in which the owner installs \emph{grants},
bounded permissions that let an app or agent act without a signature prompt for
each call. The account decides each authorization with one pure function,
rule-core's \code{rule\_eval::}\allowbreak\code{authorize}~\cite{rule-core},
linked into the account; the contract code around it, the account \emph{shell},
stores state, selects rules and checks signatures. The Perch paper,
\perchdoc{main}, states five theorems about the function; this companion gives
the definitions and all thirteen statements, and says what carries them to the
deployed account, what they trust and what they leave unproved. Theorem and
definition numbers are shared by every Perch paper.

\paragraph{Objects.} Each statement is about one or two of five objects, none of
them the compiled contract: the \emph{model}, rule-core's hand-written Lean model
of the evaluator; the \emph{translation}, the Lean program Aeneas generates from
the Rust; the \emph{account model}, an abstract account with two events and no
Rust counterpart; the \emph{host model} of Soroban authorization, whose fidelity
to soroban-env-host is tested; and Perch's \emph{canonical-encoding package}.
Aeneas gives every translated function three outcomes: success
$\mathsf{ok}(v)$, failure (a Rust panic) and divergence. A theorem that a function
returns $\mathsf{ok}$ on every input thus rules out panics and nontermination; a
refusal is a success, Rust's \code{Result<\_, Refusal>} being the value inside
$\mathsf{ok}$.

\paragraph{What ``proved'' means.} We say \emph{proved} only of statements Lean's
kernel has checked. rule-core's CI pins all \fact{exported-theorems} exported
theorems to Lean's standard axioms, \fact{standard-axioms}, and checks a ledger of
\fact{claims-total} claims: \fact{claims-proven} proven, \fact{claims-tested}
tested, \fact{claims-assumed} assumed and \fact{claims-unproven} unproven
(\S\ref{sec:trust}).

\paragraph{Reading a statement.} Each statement gives its formula, its precise
text, its hypotheses and its evidence: the Lean names, linked to their lines in
rule-core at \fact{rule-core-pin-short} (in Perch at \fact{perch-pin-short} for
Theorem~\ref{stmt:canonical-identity}), and a tag naming the object.
\statusobj{Model + Translation} means proved for the model and, under the
\lean{Extracted} namespace, for the translation. The HTML
edition\footnote{\url{https://bitswell-hq.github.io/stellar-registry__perch/perch-paper/verification/}}
adds a one-line reading and, where it helps, an instance to try.
\S\S\ref{sec:pipeline}--\ref{sec:host} state the theorems by object,
\S\ref{sec:obligations} carries them to the account, and the rest say what they
trust and leave open.
